% !TeX program = xelatex
% !BIB program = biber
% !TeX encoding = UTF-8
% ============================================================================
%  main.tex  --  THIS PRESENTATION
% ----------------------------------------------------------------------------
%  This file is the presentation. One copy of this repo = one talk.
%
%  To start a new talk: copy this whole folder, rename it, and edit
%  metadata.tex and sections/. Nothing inside needs repointing.
%
%  Build it, whichever you prefer:
%     TeXstudio / TeXworks / Overleaf   just press the green build button
%     terminal, from this folder        latexmk
%     terminal, from this folder        .\scripts\build.ps1     (or build.sh)
%
%  The three "% !TeX" lines above tell TeXstudio which engine and which
%  bibliography tool to use. Leave them at the top of the file.
%
%  Engine is XeLaTeX because Thai is switched on below. If your talk has no
%  Thai in it at all, you can change line 1 to `pdflatex` and set
%  $pdf_mode = 1 in .latexmkrc for a slightly faster build.
% ============================================================================

\documentclass[t,dvipsnames,10pt,aspectratio=169]{beamer}
%              | |         |     |
%              | |         |     +-- 16:9 widescreen (use 43 for old 4:3)
%              | |         +-------- base font size
%              | +------------------ unlocks 68 extra colour names
%              |                     (RoyalBlue, BrickRed, ForestGreen, ...)
%              +-------------------- content starts at the TOP of the slide
%
%  About that `t': without it beamer centres the body of every frame
%  vertically, so a slide with three bullets on it pushes them halfway down
%  the page and leaves a wide gap under the header. Delete the `t' to get
%  that centring back, or write \begin{frame}[c] on the one slide that wants
%  it. See the tuning block below for the rest of the vertical spacing.

% ---------------------------------------------------------------------------
%  WHERE THE SHARED FILES LIVE  --  nothing to edit here, ever
%
%  This searches upward for the folder holding theme/ and preamble/: first
%  beside this file, then one, two and three levels up. The first folder that
%  actually contains theme/beamerthemedeck.sty wins, and \deckroot is set to
%  it. Everything else in the deck is relative to \deckroot.
%
%  That search is what makes this folder copyable. All of these build as-is:
%     copy the whole repo    main.tex sits beside theme/          -> ./
%     drop it into a project one folder deeper                    -> ../
%     keep a second talk in a subfolder of this one               -> ../
%
%  No path to edit, no environment variable, no build script, no export step.
%  If nothing is found it stops immediately with a message saying so, instead
%  of failing later with a confusing "file not found".
% ---------------------------------------------------------------------------
\makeatletter
\def\deck@probe#1{%
  \ifx\deckroot\@undefined
    \IfFileExists{#1theme/beamerthemedeck.sty}{\xdef\deckroot{#1}}{}%
  \fi
}
\deck@probe{}
\deck@probe{../}
\deck@probe{../../}
\deck@probe{../../../}
\ifx\deckroot\@undefined
  \PackageError{deck}{Cannot find theme/beamerthemedeck.sty}{%
    Searched ./ ../ ../../ and ../../../ from this file.^^J%
    This deck needs the theme/ and preamble/ folders somewhere at or above
    it. Copy them next to main.tex, or copy the whole repo.}%
\fi
\edef\input@path{{\deckroot theme/}{\deckroot preamble/}}
\makeatother

% --- Theme -----------------------------------------------------------------
\usetheme{deck}
% Options:  \usetheme[noprogressbar]{deck}
%           \usetheme[nofooter]{deck}

% --- Shared setup ----------------------------------------------------------
\input{base}      % always
\input{math}      % equations          -- delete if unused
\input{code}      % code listings      -- delete if unused
\input{figures}   % TikZ / pgfplots    -- delete if unused
\input{macros}    % your own shortcuts

% --- Thai ------------------------------------------------------------------
%  On, because this template is used for Thai talks. Type Thai anywhere --
%  titles, headings, body -- with no wrapping macro. Needs XeLaTeX, which
%  line 1 and .latexmkrc already select.
%
%  Delete this line if the talk is English-only; it saves a little build time.
\input{lang-thai}

% --- Bibliography ----------------------------------------------------------
%  \cite{key} prints [1] where you write it AND the full reference, small, at
%  the bottom of that slide. The same work keeps the same number on the
%  references slide at the end (\bibliographyframe, further down).
%
%  \citequiet{key} gives the number without the bottom line.
%
%  Delete these two lines and \bibliographyframe if the talk cites nothing.
\input{bib}
\addbibresource{refs.bib}

% --- Tuning ----------------------------------------------------------------
%  Everything adjustable, in one place. All four are optional; the values
%  shown are the defaults, so uncomment a line only to change it.
%
%  Vertical spacing, if the slide body still does not sit where you want:
% \decktitlegap{1.2ex}        % gap between the frame title and the body
% \decklinespread{1.25}       % line spacing; 1.15 is tighter, see base.tex
% \setlength{\parskip}{0.5em} % space between paragraphs and list items
%
%  The per-slide reference block:
% \deckcitefootsize{\tiny}    % \scriptsize if \tiny is too small to read
% \deckcitefootrule{0.4pt}    % 0pt removes the hairline above the block
% \deckcitefootmax{4}         % warn in the log past this many on one slide
% \deckcitefootoff            % stop adding them entirely (also works inside
%                             % a single frame, where it applies to that
%                             % frame only)

% --- Who / what / when -----------------------------------------------------
% ============================================================================
%  metadata.tex  --  who is presenting, what, and when
% ----------------------------------------------------------------------------
%  The [short] versions in square brackets are what appears in the footer of
%  every slide. Keep them short or the footer will wrap.
% ============================================================================

\title[Short title]{Full Presentation Title}
\subtitle{An optional subtitle}

\author[A. Author]{Author Name}

\institute{Department \\ University}

\date{\today}

% --- Optional logos ---------------------------------------------------------
%  Point these at your own file and uncomment. A bare file name is enough for
%  anything in figures/ or theme/assets/.
%
%  Title slide, bottom right in the white band (the band is 9mm tall):
% \decklogo{\includegraphics[height=6mm]{logo-placeholder.pdf}}
%
%  Every slide, in the footer between the author and the slide number:
% \deckfooterlogo{\includegraphics[height=5mm]{logo-placeholder.pdf}}


% ============================================================================
\begin{document}

\begin{frame}[plain]
  \titlepage
\end{frame}

\begin{frame}{Outline}
  \tableofcontents
\end{frame}

% !TeX root = ../main.tex
% ============================================================================
%  Section 1 -- Introduction
% ============================================================================

\section{Introduction}

\begin{frame}{Motivation}
  \begin{itemize}
    \item First point, resting on earlier work \cite{example2024}
    \item Second point
      \begin{itemize}
        \item A supporting detail
      \end{itemize}
    \item Third point, and two sources at once \cite{example2024,sample2023}
  \end{itemize}

  % The two \cite commands above put [1] and [2] in the text, the matching
  % full references at the bottom of this slide, and both works on the
  % references slide at the end. Note that citing [1] twice on one slide
  % still gives one line at the bottom. Delete them with the placeholder text.
\end{frame}

\begin{frame}{Contributions}
  \begin{block}{What this work does}
    A short statement of the contribution.
  \end{block}
\end{frame}

% ============================================================================
%  DELETE THIS SLIDE -- it is here so a fresh copy proves Thai is working
%  ลบสไลด์นี้ทิ้งได้ มีไว้เพื่อยืนยันว่าภาษาไทยใช้งานได้ตั้งแต่เริ่ม
% ============================================================================
\begin{frame}{Thai check / ตรวจสอบภาษาไทย}
  ถ้าคุณอ่านบรรทัดนี้ออก แปลว่าฟอนต์ภาษาไทยทำงานถูกต้องแล้ว
  พิมพ์ภาษาไทยได้ทุกที่ ทั้งหัวเรื่องและเนื้อหา โดยไม่ต้องครอบคำสั่งใด ๆ

  \vspace{0.4em}
  If you can read the line above, Thai is working. Type it anywhere ---
  titles included --- with no wrapping macro. Bold carries across:
  \textbf{Results / ผลลัพธ์}.
\end{frame}

% !TeX root = ../main.tex
% ============================================================================
%  Section 2 -- Method
% ============================================================================

\section{Method}

\begin{frame}{Approach}
  \twocol{%
    \begin{itemize}
      \item Step one
      \item Step two
    \end{itemize}
  }{%
    Right-hand column: a figure, a table, or more text.
  }
\end{frame}


% --- References ------------------------------------------------------------
%  Every work cited anywhere in the talk, numbered as it was cited. Splits
%  across several slides on its own if the list is long.
\bibliographyframe

% --- Appendix --------------------------------------------------------------
%  Slides after \appendix are not counted in the "N / total" footer.
\appendix
% !TeX root = ../main.tex
% ============================================================================
%  Appendix -- backup slides
% ----------------------------------------------------------------------------
%  These come after \appendix in main.tex, so they do not count towards the
%  "N / total" number in the footer. Useful for questions after the talk.
% ============================================================================

\begin{frame}{Backup: additional detail}
  Extra material kept in reserve for questions.
\end{frame}


\end{document}
